\documentclass[12pt]{extarticle}

\usepackage[a4paper,margin=1in]{geometry}
\usepackage{amsmath,amssymb,amsfonts}
\usepackage{enumitem}
\usepackage{hyperref}

\setlength{\parskip}{0.55em}
\setlength{\parindent}{0pt}

\title{Code Overview}
\author{}
\date{}

\begin{document}
\maketitle

\section{Implementation Rationale}

This section describes the currently implemented codebase as the practical realization of the dataset methodology proposed in this thesis. The core idea is that humor data should not be treated as a flat collection of isolated jokes. Instead, it should be transformed into a layered resource in which each joke is connected to semantic representations, weak constraints, and multiple references. Such a representation is better aligned with creative generation, where a single prompt may admit many acceptable outputs and where exact lexical overlap is often a poor proxy for semantic relatedness.

The repository is therefore not merely an engineering scaffold. It operationalizes a research assumption: before one can study reasoning-based humor generation, one needs a corpus that is both large and semantically structured. The implemented pipeline moves from raw text collection to progressively richer annotations, and each stage materializes a separate dataset. Methodologically, this is useful because it makes the process reproducible, allows each representation layer to be inspected independently, and avoids collapsing the whole construction problem into a single brittle model call.

A second motivation behind the current design is epistemic caution. In humor research, many desirable labels are difficult to define sharply and even harder to annotate consistently at scale. It is therefore advantageous to build supervision gradually. Dense embeddings, sparse keywords, and retrieved references are not final ground truth in the strict sense, but they are informative intermediate objects that can support later training and analysis. The implemented code should thus be read as a pipeline for constructing increasingly informative approximations to the latent structure of humorous texts.

This staged methodology also matches the broader argument of the thesis. The long-term objective is not simply to store jokes, but to organize them in a way that supports reasoning-oriented generation and learning without relying on a single verifier. For such a setting, the corpus must expose relations between jokes: topical proximity, shared constraints, and families of plausible reference answers. The existing code already moves in this direction by turning a raw corpus into a set of linked semantic views.

\section{Dataset Construction Strategy}

\textbf{Corpus-first design.} The implemented pipeline begins with large externally available joke corpora rather than with manually authored prompts. This reflects a methodological trade-off: if the objective is to construct a resource with diverse humor mechanisms, broad topical coverage, and many potential reference answers, then naturally occurring joke collections are preferable to small curated prompt-response datasets. They contain recurring topics, repeated lexical cues, and stylistic variation that later make multi-reference supervision possible.

\textbf{Source selection and complementarity.} The current code uses two fixed sources: the \emph{Short Jokes} corpus and the \emph{rJokes} corpus. Their combination is methodologically attractive because it balances short one-liners with more heterogeneous community-style jokes. The first source contributes brevity, lexical compactness, and many formulaic punchline patterns. The second contributes broader topical diversity, looser discourse structure, and a wider range of community-generated humor forms. Using both sources is therefore a way of reducing dependence on a single stylistic regime.

\textbf{Reproducibility and frozen acquisition.} The sources are downloaded from pinned GitHub commits rather than mutable endpoints, so the acquisition process remains stable across runs. This detail matters methodologically because later analysis of retrieval quality, prompt construction, or training behavior should not be confounded by silent changes in the underlying corpus. Freezing acquisition reduces dataset drift and supports later comparison and ablation studies.

\textbf{Canonicalization and provenance preservation.} After acquisition, the raw files are normalized into a shared parquet representation with a unified global identifier, while source metadata are preserved. This reflects the methodological principle that heterogeneous corpora should be merged without erasing provenance. Retaining source information makes it possible to inspect domain effects, identify source-specific artifacts, or later perform filtering and balancing, while the global identifier becomes the stable key for all downstream annotations.

\textbf{Minimal normalization and noise tolerance.} The implemented collection stage performs only lightweight cleaning: empty records are removed, identifiers are normalized, and all examples are rewritten into a common tabular structure. This restraint is methodologically deliberate. Over-aggressive text normalization can destroy exactly those surface irregularities that are relevant for humor, including orthographic play, punchline timing, and deliberately awkward formulations. The current code therefore prefers a permissive collection strategy in which noise is preserved unless it directly prevents downstream processing.

\textbf{A staged view of semantic structure.} The codebase does not jump directly from raw jokes to prompts or references. Instead, it inserts intermediate representational layers. This is motivated by the fact that humor is semantically dense but lexically unstable: two jokes may be close in premise or mechanism even if they share little surface vocabulary. A staged pipeline therefore offers a more reliable route from text to weak supervision than exact matching would.

\textbf{Expected distributional properties.} From a methodological point of view, the collected corpus should be expected to be uneven. Certain topics, lexical triggers, and joke forms occur frequently, while others are rare. This is not merely an inconvenience; it is part of the object of study. A prompt-oriented humor dataset naturally has a heavy-tailed structure in which some weak constraints map to many jokes and others map to very few. The current implementation preserves this structure rather than artificially smoothing it out, because later stages such as retrieval and reference construction depend precisely on such non-uniform semantic neighborhoods.

\section{Semantic Enrichment and Reference Construction}

\textbf{Dense embeddings as a continuous semantic layer.} The first enrichment stage maps each joke to an embedding vector. Embeddings provide a continuous semantic space in which related jokes can become neighbors even when their wording differs. This is especially useful for humor, where paraphrase, indirection, and lexical novelty are common. By materializing embeddings as a standalone dataset, the pipeline creates the geometric backbone for later retrieval operations.

\textbf{Why embeddings come before prompt construction.} In principle, one could attempt to derive prompts directly from raw jokes with a language model. The current implementation does not begin there. Instead, it first constructs a geometric semantic layer. Methodologically, this provides a more stable substrate on which later approximations can operate. A dense representation makes it possible to compare jokes, keywords, and weak prompts within a common space, which is valuable when exact symbolic mappings are not yet available.

\textbf{Keywords as an interpretable bottleneck.} The next stage extracts a small set of keywords for each joke. This step should be understood not as naive keyword mining, but as semantic compression. The aim is to produce a sparse and interpretable representation that captures salient constraints without preserving the whole text. In the implemented system, lexical candidates are generated from the joke itself, embedded, and ranked against the embedding of the full joke. Maximal marginal relevance balances relevance and diversity, so the final keywords function as compact semantic handles that remain human-readable while still being grounded in the embedding space.

\textbf{The role of diversity in keyword selection.} Diversity is an important methodological consideration here. If keyword extraction simply returned the most frequent or most semantically similar terms, the output would often collapse into near-duplicates and fail to express the constraint set of the joke. The use of diversity-aware selection acknowledges that a useful weak prompt should ideally reflect several facets of a joke at once, for example topic, object, setting, or mechanism. Even though the current implementation is lightweight, it already treats keyword extraction as a constrained summarization problem rather than as a purely lexical one.

\textbf{Retrieval as practical prompt inversion.} The broader thesis proposes converting jokes into prompts so that many jokes can be grouped under shared weak constraints. The current implementation realizes a practical approximation of this idea. Instead of asking a language model to infer prompts for every joke, it forms weak prompt-like queries from keyword combinations and retrieves semantically similar jokes in embedding space. Conceptually, this is an inversion-by-retrieval strategy: a joke is compressed into salient cues, these cues are expanded into weak prompts, and those prompts are used to recover a neighborhood of plausible references.

This approximation is methodologically defensible. Direct prompt extraction is flexible, but it is also expensive and may introduce instability or hallucinated constraints. The implemented retrieval-based route sacrifices some expressivity in exchange for reproducibility and control. It still preserves the key structural idea needed for later training: the same weak semantic seed can point to multiple acceptable jokes. Thus, the current code already instantiates a multi-reference view of humor, even if it does so through keyword-guided retrieval rather than full prompt induction.

\textbf{From keyword groups to weak constraints.} A particularly important feature of the implementation is that it does not restrict itself to a single keyword view of each joke. Instead, it forms several keyword groups and turns them into prompt-like queries. Methodologically, this matters because humorous texts are often multiply interpretable. One joke can be indexed through a topical cue, a relation between concepts, or a more specific object or phrase. Generating several weak constraints from the same joke is therefore a way of approximating the fact that a single humorous text may belong to several semantic families at once.

\textbf{Multi-reference structure as the central target.} The ultimate goal of these transformations is not retrieval for its own sake. It is the construction of a dataset in which one weak semantic seed is linked to multiple plausible joke references. This is the key property needed for the broader thesis, because it shifts the supervision problem away from a single correct answer and toward a partially observed set of good answers. The current code reaches this structure operationally by storing, for each identifier and keyword group, a list of retrieved references with associated similarity scores.

\textbf{Indexing and scalability.} The use of FAISS for reference retrieval is also part of the methodological design rather than a purely technical convenience. A multi-reference dataset is only useful if it can be constructed over hundreds of thousands of examples, which makes naive all-pairs comparison impractical. Approximate nearest-neighbor search therefore provides a scalable approximation to semantic grouping.

\textbf{Methodological limitations of retrieval-based inversion.} At the same time, the current implementation should not be mistaken for a perfect realization of prompt induction. Retrieval neighborhoods can reflect superficial topical similarity without fully preserving the underlying comedic mechanism. Keyword-based weak constraints may underspecify the joke, overspecify it, or privilege objects over relations. In addition, nearest-neighbor retrieval can return the original joke or jokes that are semantically adjacent but not truly interchangeable as references. These limitations do not invalidate the approach, but they clarify what kind of supervision the present pipeline actually constructs: a pragmatic approximation to a multi-reference prompt dataset, not a final canonical form of it.

\textbf{Why the approximation is still useful.} Despite these limitations, the current design remains valuable for research. It already enforces the shift from isolated examples to structured neighborhoods; it already separates dense and sparse representations; and it already exposes candidate reference sets that can be inspected, filtered, or rescored later. In other words, the implemented pipeline does not solve the full prompt construction problem, but it creates the data substrate on top of which a more explicit prompt-centric method can be added.

\section{Downstream Experimental Utility}

\textbf{Evaluation as a separate but necessary layer.} The evaluation pipeline is not itself a data collection stage, yet it plays an important methodological role in the overall system. It takes externally generated candidate jokes, builds pairwise comparisons within the same prompt, and aggregates judgments into leaderboard-style summaries. This reflects a broader assumption already discussed in the thesis: humor quality is difficult to verify directly, so practical evaluation often relies on comparative judgment rather than absolute scoring.

\textbf{Pairwise comparison as a natural evaluation format.} This design is especially appropriate for humor. Absolute scores for funniness are unstable and context-dependent, whereas pairwise judgments ask a weaker and often more reliable question: which of two responses is better for the same prompt? By materializing pairwise outputs and then aggregating them into model-level summaries, the current code follows a methodology that is better aligned with the ambiguity of humor evaluation than single scalar ratings would be.

\textbf{Leaderboard construction as an experimental interface.} The final leaderboard should therefore be interpreted not as a benchmark in the strict sense, but as an interface for model comparison within a controlled experimental setting. Once the data construction pipeline produces references and weak constraints, the evaluation pipeline makes it possible to test whether those structures are actually useful for generation. In this sense, the repository closes a methodological loop: the same project that builds structured supervision also provides a way to assess whether that supervision leads to better outputs.

\textbf{Materialized intermediate datasets.} Another notable property of the current implementation is that every stage writes its outputs as parquet shards. This supports resumability and large-scale processing, but it also has methodological value. It makes each transformation explicit: jokes, embeddings, keywords, and references are all inspectable as separate datasets. For a thesis project, this is useful because it allows failure analysis at the level of representation rather than only at the final generation stage.

\textbf{Resumability and large-scale iteration.} The sharded structure of the outputs also supports a practical research workflow in which collection and enrichment can be rerun incrementally. Although this may seem like an engineering detail, it has methodological consequences: when a dataset is built in stages, one often wants to improve a single layer without discarding all previous computation. The current code supports this style of iterative experimentation, which is particularly useful when embedding models, retrieval parameters, or filtering criteria are still under active investigation.

\textbf{Current scope of the implementation.} Finally, it is important to distinguish between the full conceptual pipeline and the subset that is already implemented. The repository already supports reproducible source acquisition, canonicalization, dense semantic encoding, interpretable keyword extraction, retrieval-based multi-reference construction, and pairwise evaluation. However, it does not yet fully implement the model-based prompt extraction and normalization procedure envisioned in the broader methodological plan.

The current code should therefore be interpreted as an intermediate but coherent realization of the thesis direction. It already embodies the main design principles of weak constraints, multiple references, and layered semantic structure, while leaving room for a more explicit prompt-centric formulation in future work. From a methodological perspective, this is a meaningful stage of progress: the project has already moved beyond a raw joke corpus and toward a structured humor dataset that can support retrieval, analysis, and future learning objectives.

\end{document}
